% Canonical manuscript source. The older Pandoc Markdown source was archived in
% ../archive/2026-04-03-markdown-source/gut_application_paper.md.
% Options for packages loaded elsewhere
\PassOptionsToPackage{unicode}{hyperref}
\PassOptionsToPackage{hyphens}{url}
\PassOptionsToPackage{dvipsnames,svgnames,x11names}{xcolor}
\documentclass[
  11pt,
]{article}
\usepackage{xcolor}
\usepackage[margin=1in]{geometry}
\usepackage{amsmath,amssymb}
\setcounter{secnumdepth}{-\maxdimen} % remove section numbering
\usepackage{iftex}
\ifPDFTeX
  \usepackage[T1]{fontenc}
  \usepackage[utf8]{inputenc}
  \usepackage{textcomp} % provide euro and other symbols
\else % if luatex or xetex
  \usepackage{unicode-math} % this also loads fontspec
  \defaultfontfeatures{Scale=MatchLowercase}
  \defaultfontfeatures[\rmfamily]{Ligatures=TeX,Scale=1}
\fi
\usepackage{lmodern}
\ifPDFTeX\else
  % xetex/luatex font selection
\fi
% Use upquote if available, for straight quotes in verbatim environments
\IfFileExists{upquote.sty}{\usepackage{upquote}}{}
\IfFileExists{microtype.sty}{% use microtype if available
  \usepackage[]{microtype}
  \UseMicrotypeSet[protrusion]{basicmath} % disable protrusion for tt fonts
}{}
\usepackage{setspace}
\makeatletter
\@ifundefined{KOMAClassName}{% if non-KOMA class
  \IfFileExists{parskip.sty}{%
    \usepackage{parskip}
  }{% else
    \setlength{\parindent}{0pt}
    \setlength{\parskip}{6pt plus 2pt minus 1pt}}
}{% if KOMA class
  \KOMAoptions{parskip=half}}
\makeatother
\usepackage{longtable,booktabs,array}
\usepackage{calc} % for calculating minipage widths
% Correct order of tables after \paragraph or \subparagraph
\usepackage{etoolbox}
\makeatletter
\patchcmd\longtable{\par}{\if@noskipsec\mbox{}\fi\par}{}{}
\makeatother
% Allow footnotes in longtable head/foot
\IfFileExists{footnotehyper.sty}{\usepackage{footnotehyper}}{\usepackage{footnote}}
\makesavenoteenv{longtable}
\usepackage{graphicx}
\makeatletter
\newsavebox\pandoc@box
\newcommand*\pandocbounded[1]{% scales image to fit in text height/width
  \sbox\pandoc@box{#1}%
  \Gscale@div\@tempa{\textheight}{\dimexpr\ht\pandoc@box+\dp\pandoc@box\relax}%
  \Gscale@div\@tempb{\linewidth}{\wd\pandoc@box}%
  \ifdim\@tempb\p@<\@tempa\p@\let\@tempa\@tempb\fi% select the smaller of both
  \ifdim\@tempa\p@<\p@\scalebox{\@tempa}{\usebox\pandoc@box}%
  \else\usebox{\pandoc@box}%
  \fi%
}
% Set default figure placement to htbp
\def\fps@figure{htbp}
\makeatother
% definitions for citeproc citations
\NewDocumentCommand\citeproctext{}{}
\NewDocumentCommand\citeproc{mm}{%
  \begingroup\def\citeproctext{#2}\cite{#1}\endgroup}
\makeatletter
 % allow citations to break across lines
 \let\@cite@ofmt\@firstofone
 % avoid brackets around text for \cite:
 \def\@biblabel#1{}
 \def\@cite#1#2{{#1\if@tempswa , #2\fi}}
\makeatother
\newlength{\cslhangindent}
\setlength{\cslhangindent}{1.5em}
\newlength{\csllabelwidth}
\setlength{\csllabelwidth}{3em}
\newenvironment{CSLReferences}[2] % #1 hanging-indent, #2 entry-spacing
 {\begin{list}{}{%
  \setlength{\itemindent}{0pt}
  \setlength{\leftmargin}{0pt}
  \setlength{\parsep}{0pt}
  % turn on hanging indent if param 1 is 1
  \ifodd #1
   \setlength{\leftmargin}{\cslhangindent}
   \setlength{\itemindent}{-1\cslhangindent}
  \fi
  % set entry spacing
  \setlength{\itemsep}{#2\baselineskip}}}
 {\end{list}}
\usepackage{calc}
\newcommand{\CSLBlock}[1]{\hfill\break\parbox[t]{\linewidth}{\strut\ignorespaces#1\strut}}
\newcommand{\CSLLeftMargin}[1]{\parbox[t]{\csllabelwidth}{\strut#1\strut}}
\newcommand{\CSLRightInline}[1]{\parbox[t]{\linewidth - \csllabelwidth}{\strut#1\strut}}
\newcommand{\CSLIndent}[1]{\hspace{\cslhangindent}#1}
\setlength{\emergencystretch}{3em} % prevent overfull lines
\providecommand{\tightlist}{%
  \setlength{\itemsep}{0pt}\setlength{\parskip}{0pt}}
\usepackage{booktabs}
\usepackage{longtable}
\usepackage{array}
\usepackage{tabularx}
\usepackage{float}
\usepackage{titling}
\usepackage{caption}
\captionsetup{font=small,labelfont=bf}
\newcolumntype{Y}{>{\raggedright\arraybackslash}X}
\newcommand{\taxon}[1]{\emph{#1}}
\newcommand{\dcstsep}{/\allowbreak}
\setlength{\droptitle}{-2.6em}
\pretitle{\begin{center}\Large}
\posttitle{\par\end{center}\vspace{0.55em}}
\preauthor{\begin{center}\normalsize}
\postauthor{\par\end{center}\vspace{-0.5em}}
\predate{}
\postdate{}
\usepackage{bookmark}
\IfFileExists{xurl.sty}{\usepackage{xurl}}{} % add URL line breaks if available
\urlstyle{same}
\hypersetup{
  pdftitle={Interpretable Dominant-Taxon States for Characterizing IBD and Related Gut Microbiome Dysbiosis},
  colorlinks=true,
  linkcolor={blue},
  filecolor={Maroon},
  citecolor={Blue},
  urlcolor={blue},
  pdfcreator={LaTeX via pandoc}}
\newcommand{\paperaffiliation}{Center for Advanced Microbiome Research and Innovation (CAMRI), Institute for Genome Sciences, Department of Microbiology and Immunology, University of Maryland School of Medicine, Baltimore, Maryland, USA.}
\newcommand{\paperfundingnote}{Research reported in this publication was supported by grant number INV-048956 from the Gates Foundation.}
\newcommand{\blfootnote}[1]{%
  \begingroup
  \renewcommand{\thefootnote}{}\footnote{#1}%
  \addtocounter{footnote}{-1}%
  \endgroup
}

\title{Interpretable Dominant-Taxon States for Characterizing IBD and
Related Gut Microbiome Dysbiosis}
\author{\parbox{0.9\textwidth}{\centering Pawel Gajer\thanks{\paperaffiliation} and Jacques Ravel\protect\footnotemark[1]}}
\date{}
\newif\ifshowbuildstamp
\showbuildstampfalse
\providecommand{\manuscriptversion}{unversioned}
\providecommand{\manuscriptbuildnumber}{NA}
\providecommand{\manuscriptcommit}{unknown}
\providecommand{\manuscriptbuilddatetime}{unknown}
\InputIfFileExists{../build/manuscript_build_info.tex}{}{}

\begin{document}
\maketitle
\blfootnote{\paperfundingnote}
\ifshowbuildstamp
\begin{center}
\footnotesize
% Version \texttt{\manuscriptversion} \quad | \quad Build \texttt{\manuscriptbuildnumber} \quad | \quad Commit \texttt{\manuscriptcommit}\\
Built \manuscriptbuilddatetime
\end{center}
\vspace{0.35em}
\fi

\setstretch{1.15}
\section{Abstract}\label{abstract}

Inflammatory bowel disease (IBD) has repeatedly been associated with
heterogeneous gut microbiome dysbiosis, including loss or redistribution
of obligate anaerobic commensals and expansion of facultative or
inflammation-associated taxa. The practical challenge addressed here is
not whether these broad taxonomic shifts occur, but whether they can be
summarized as sample-level ecological states that remain interpretable
across large cohorts. Gut enterotypes provide one such summary, but they
are clustering-defined constructs whose labels depend on choices about
data transformation, distance metric, clustering model, and the number of
clusters. We therefore applied dominant-community state types (dCSTs), a
deterministic rank-based description of retained-feature dominance, to
30,290 quality-controlled American Gut Project stool samples from a
34,711-sample SILVA gut abundance table. Self-reported IBD showed the
strongest phenotype-level association with the dominance-state
hierarchy. The central IBD result was that the familiar broad
\taxon{Bacteroides}-associated dysbiosis pattern resolved into an
ordered \taxon{Bacteroides}/Lachnospiraceae/\taxon{Alistipes}
dominance-lineage, providing a more specific sample-level description of
IBD-associated ecological context rather than implying that the named
taxa are newly discovered IBD markers. Bayesian models also provided stabilized estimates for sparse
\taxon{Morganella} and \taxon{Proteus} states, which are treated as
exploratory enterobacterial signals. The IBD dominance-state summary was
robust to a stricter control definition and to coarse antibiotic-status
sensitivity analyses. External validation separated two forms of
reproducibility: sample-level Halfvarson 2017 analyses supported both
rebuilt-cohort and direct AGP-derived label-transfer signals,
PRJEB84421 provided narrower rebuilt inflammatory support, and pooled
HMP2 retained a weaker depth-2 transferred IBD signal that persisted in a
first-sample-per-subject sensitivity analysis. IBS and autoimmune disease
showed weaker but directionally related dominance patterns. These results
support dCSTs as an interpretable framework for organizing established
IBD-associated dysbiosis into deterministic dominance-lineage states,
while broader gut-outcome applications and universal label portability
remain exploratory.

\textbf{Keywords:} gut microbiome; inflammatory bowel disease;
dysbiosis; enterotypes; dominant-community state types; Bayesian
modeling

\section{Introduction}\label{introduction}

For common gut disorders, clinical phenotypes do not map neatly onto one
standalone molecular or laboratory test. IBD diagnosis integrates
clinical presentation with biochemical, stool, \mbox{endoscopic},
imaging, and histologic evidence rather than relying on one reference standard
(\citeproc{ref-Maaser2019_ECCO_ESGAR_IBD_diagnostics}{Maaser et al.
2019}). IBS is even more symptom-defined, with no validated diagnostic
biomarker and a guideline-supported positive diagnostic strategy
(\citeproc{ref-Lacy2021_ACG_IBS}{Lacy et al. 2021}). This creates a
substantial interpretive challenge for microbiome studies:
disease-associated community shifts are repeatedly observed, but they are
heterogeneous and difficult to translate into sample-level ecological
labels that can be compared across cohorts.

In IBD, the existence of microbiome dysbiosis is no longer the central
question. Landmark studies have shown that IBD is not characterized by
one universal microbial signature, but often involves reduced diversity,
depletion or redistribution of butyrate-associated obligate anaerobes,
and expansion of facultative or inflammation-associated taxa
(\citeproc{ref-Frank2007_IBD}{Frank et al. 2007};
\citeproc{ref-Gevers2014_CD}{Gevers et al. 2014};
\citeproc{ref-LloydPrice2019_IBD}{Lloyd-Price et al. 2019};
\citeproc{ref-Sokol2008_Fprausnitzii}{Sokol et al. 2008};
\citeproc{ref-Manichanh2006_CD}{Manichanh et al. 2006}). These
observations motivate a practical representation question rather than a
new diagnostic-test claim: can established but heterogeneous
IBD-associated microbiome shifts be summarized into transparent
sample-level dominance states that can be compared across cohorts?

Gut enterotypes were introduced to provide such labels by compressing
high-dimensional gut abundance profiles into a small number of
clustering-defined configurations
(\citeproc{ref-Arumugam2011_enterotypes}{Arumugam et al. 2011}). That
abstraction can be informative for broad community structure, but it has
remained controversial precisely because the labels are clustering
constructs: assignments depend on data transformation, distance metric,
initialization, model class, and the chosen number of clusters
(\citeproc{ref-Koren2013_enterotypes}{Koren et al. 2013};
\citeproc{ref-Knights2014_rethinking_enterotypes}{Knights et al. 2014};
\citeproc{ref-Costea2018_enterotypes}{Costea et al. 2018}). Those
choices can make it hard to determine what a label means biologically
for a particular sample or whether labels from different analyses are
directly comparable.

We use dominant-community state types (dCSTs) as a deterministic
rank-based framework for naming sample-level gut ecological states. A
dCST is defined by the within-sample ordering of retained taxa. At depth
1, the state is assigned by the dominant retained taxon; at depth 2, it
is refined by the ordered pair of the two most abundant retained taxa.
Because assignment depends only on within-sample rank order, the same
label is recovered from counts, proportions, or any sample-wide
rescaling that preserves taxon order within a sample. We use
\emph{dominance-lineage} to denote the ordered path defined by a
depth-\(\ell\) dCST label as the hierarchy is refined across ranks. The
term refers only to within-sample dominance order and does not imply
phylogenetic relatedness. The dCST construction is therefore not intended
to replace mechanistic microbiome modeling. Instead, it provides a
transparent way to ask whether IBD-associated dysbiosis can be summarized
as dominance states that remain interpretable after covariate adjustment,
hierarchical refinement, Bayesian sparse-state follow-up, and external
validation. The construction is illustrated in Figure 1.

\begin{figure}[H]
\centering
\includegraphics[width=0.96\linewidth,height=\textheight,keepaspectratio,alt={Conceptual schematic of dCST construction. Panel A shows a toy gut abundance table. Panel B groups samples into dominance sample sets by their top-ranked retained taxon. Panel C contrasts the pure policy, which sends small dominance sample sets to a residual bucket, with the absorb policy, which merges them into retained named states. Panel D shows that a depth-2 label is an ordered dCST dominance-lineage such as Bacteroides/Faecalibacterium.}]{../assets/figures/FIGURE_1_dcst_conceptual_schematic.png}
\caption{Conceptual schematic of dCST construction. Panel A shows a toy
gut abundance table, and panel B groups samples into dominance sample
sets (DSSs) according to their top-ranked retained taxon. Panel C
contrasts the pure-policy view, in which low-frequency DSSs collapse
into a residual bucket, with the absorb-policy view used here, in which
low-frequency patterns are merged into retained named states. Panel D
shows that a depth-2 label is an ordered dCST dominance-lineage, such as
\taxon{Bacteroides}/\taxon{Faecalibacterium}, rather than a fitted
cluster or a phylogenetic relationship.}
\end{figure}

This application therefore centers on IBD, with IBS and autoimmune
disease included as related gut or immune-associated outcomes because
literature suggests that shifts in dominant taxa or dominant taxon pairs
can mark altered ecological states in those settings
(\citeproc{ref-Tap2017_IBS}{Tap et al. 2017};
\citeproc{ref-Pittayanon2019_IBS}{Pittayanon et al. 2019};
\citeproc{ref-Scher2013_Prevotella_RA}{Scher et al. 2013};
\citeproc{ref-Vatanen2018_T1D}{Vatanen et al. 2018}). We use
``inflammatory'' as a disease-contextual and literature-grounded
descriptor rather than as evidence that host inflammatory biomarkers were
measured in the AGP discovery cohort. In this study we use the absorb
low-frequency policy, which merges very small dominance patterns into
larger named states rather than pooling them into one catch-all rare
category. Pure-policy dCST results are used only as a baseline
comparison in the supplement rather than as a second main-text results
analysis.

Figure 2 previews how the dCST framework organizes the IBD-centered
result: self-reported IBD had the strongest depth-resolved association
with the dominance-state hierarchy, with a broad \taxon{Bacteroides}
state resolving into more specific \taxon{Bacteroides}-centered
dominance-lineages and the strongest external support coming from
IBD-focused validation cohorts.

\begin{figure}[H]
\centering
\includegraphics[width=0.98\linewidth,height=\textheight,keepaspectratio,alt={IBD overview figure. Panel A shows that Cramer's V increases most strongly across depths 1 to 4 for IBD compared with autoimmune disease and IBS. Panel B shows frequentist and Bayesian adjusted odds ratios for representative common IBD dominance-lineages and rare enterobacterial states. Panel C compares sample-level rebuilt-cohort external validation with direct transfer of AGP-derived absorb dCST labels.}]{../assets/figures/FIGURE_2_ibd_results_overview.png}
\caption{Overview of IBD-associated dominance-state structure. Panel A
shows that IBD had the strongest depth-resolved absorb dCST association
among the main gut or immune-associated phenotype analyses. Panel B
summarizes representative common IBD dominance-lineages and sparse enterobacterial
states using frequentist and Bayesian adjusted odds ratios. Panel C
summarizes sample-level external validation: rebuilt-cohort analyses test whether
disease-associated dominance structure is rediscovered within each
cohort, whereas direct transfer of AGP-derived absorb dCST labels tests
whether discovery-cohort labels remain informative without rebuilding the
hierarchy. Subject-level validation sensitivity is reported in the text
and supplement. The dashed line marks q = 0.05.}
\end{figure}

\section{Results}\label{cohort-overview}

\subsection{Outcome-wide discovery overview}\label{outcome-wide-discovery-overview}

The discovery analysis asked whether interpretable dominance states
could organize disease-associated gut ecology in 30,290
quality-controlled AGP stool samples. The absorb dCST hierarchy
contained 40 named states at depth 1, 153 at depth 2, 479 at depth 3,
and 744 at depth 4. The largest depth-1 states were
\taxon{Bacteroides} (10,599 samples),
\taxon{Escherichia-Shigella} (6,913), \taxon{Prevotella\_9} (1,926),
Enterobacteriaceae (1,596), and \taxon{Faecalibacterium} (1,405). The
twelve most common depth-1 absorb dCST dominance-lineages covered 87.3\% of the cohort
(Supplementary Table S5), showing that a limited number of dominant
ecological configurations accounted for most samples even after
low-frequency patterns were absorbed into larger named states. Some
retained-feature labels also carried low-resolution or potential non-gut
annotations, including
\emph{Unassigned}, \taxon{Pseudomonas}, \taxon{Staphylococcus}, and, at
deeper levels, \emph{Eukaryota}/\emph{Unassigned}/Bacteria/Mitochondria
(Supplementary Tables S2 and S5). We retained these labels in the
unsupervised hierarchy for transparent accounting, but they are not used
as primary biological evidence for the IBD interpretation.

The depthwise phenotype screen showed that IBD provided the clearest
outcome-level association with the dominance-state hierarchy. Across depths 1 to 4, Cramer's V increased from
0.11 to 0.35 in IBD, compared with 0.11 to 0.28 in autoimmune disease
and 0.08 to 0.19 in IBS. Because IBD combined the largest depthwise
growth in effect size, the clearest dominance-lineage-level
interpretability, and the strongest external validation, we focus the
detailed follow-up on IBD. IBS and autoimmune disease are retained as
related secondary analyses, while additional non-IBD examples are
summarized in the supplement (Supplementary Table S6). The absorb
low-frequency policy most affected deeper levels of the hierarchy and is
summarized separately in Supplementary Table S2.

\subsection{IBD-associated dominance-lineage structure}\label{ibd-associated-dominance-lineage-structure}


Self-reported IBD provides the most clinically grounded inflammatory
phenotype in the AGP screen, even though the discovery metadata do not
separate treatment-naive Crohn disease, ulcerative colitis, disease
activity, and medication exposure.

At the label level, the primary adjusted depth-1 screen showed
enrichment of \taxon{Bacteroides} (adjusted OR 1.54, q =
8.1e-10), \taxon{Bifidobacterium} (adjusted OR 3.38, q =
2.2e-04), Enterobacterales (adjusted OR 3.23, q = 0.0016), and
Lachnospiraceae (adjusted OR 1.99, q = 0.0030), together with
depletion of \taxon{Prevotella\_9} (adjusted OR 0.44, q =
2.2e-04). The Bayesian layer recovered the same common-state picture,
but it also highlighted very rare enterobacterial states that were too
small for the primary adjusted frequentist screen. In particular,
\taxon{Morganella} showed a Bayesian adjusted OR of 18.99
(95\% credible interval 10.15--35.53) and \taxon{Proteus} showed a
Bayesian adjusted OR of 3.40 (95\% credible interval 1.40--8.27), both
with near-complete posterior support for enrichment in IBD. These
rare-state examples were small: \taxon{Morganella} contained 22 IBD
samples and 34 controls, and \taxon{Proteus} contained 6 IBD samples and
45 controls. We therefore treat them as exploratory facultative-anaerobe
or Enterobacterales-linked rare-state signals rather than as stand-alone
biomarkers or as evidence of disease severity.

Hierarchical refinement organized the broad IBD-associated pattern into
more specific dominance-lineage summaries. At depth 2,
\taxon{Bacteroides}\dcstsep Lachnospiraceae reached adjusted
OR 3.21 (q = 2.6e-33), \taxon{Bacteroides}\dcstsep[\emph{Eubacterium}]
\emph{eligens} group reached adjusted OR 22.50 (q = 5.0e-24), and
\taxon{Bacteroides}\dcstsep Enterobacteriaceae reached adjusted OR 2.73
(q = 2.2e-04). At depth 3, the strongest dominance-lineage was
\taxon{Bacteroides}\dcstsep Lachnospiraceae\dcstsep\taxon{Alistipes}
(adjusted OR 8.56, q = 1.5e-28), followed by
\taxon{Bacteroides}\dcstsep\emph{[Eubacterium]\_eligens\_group}\dcstsep
Lachnospiraceae (adjusted OR 22.50, q = 5.0e-24) and
\taxon{Bacteroides}\dcstsep\taxon{Alistipes}\dcstsep Lachnospiraceae
(adjusted OR 5.95, q = 6.1e-15). At depth 4, the best supported dominance-lineage
remained
\taxon{Bacteroides}\dcstsep Lachnospiraceae\dcstsep\taxon{Alistipes}\dcstsep\taxon{Faecalibacterium}
(adjusted OR 5.53, q = 1.1e-07). Together, these deeper states indicate
that the broad depth-1 \taxon{Bacteroides} state contains multiple
IBD-associated dominance-lineages. The absorb policy preserves these
dominance-lineages as named states rather than collapsing them into a
generic rare-state category. A dominance-lineage label records rank context rather than a claim
that every named component is enriched in IBD. For example, the depth-4
\taxon{Bacteroides}\dcstsep Lachnospiraceae\dcstsep\taxon{Alistipes}\dcstsep\taxon{Faecalibacterium}
state remains a \taxon{Bacteroides}-first retained-feature context;
\taxon{Faecalibacterium} is a lower-ranked component within that label
(Supplementary Table S8). This distinction is important because the broader
IBD literature often reports depletion of \emph{F. prausnitzii} or
butyrate-associated commensals.

The IBD-associated dominance-state summary was also robust to a stricter
control definition. After
removing non-IBD controls who carried any of the other AGP phenotype
strings analyzed in this study (IBS, autoimmune disease, acid reflux, or
seasonal allergies), the analysis retained all 1,166 IBD cases and 16,175
cleaner controls. The broad \taxon{Bacteroides} signal remained enriched
(adjusted OR 1.59, q = 1.1e-10), as did the
\taxon{Bacteroides}\dcstsep Lachnospiraceae dominance-lineage at depth 2
(adjusted OR 3.15, q = 1.5e-28) and the
\taxon{Bacteroides}\dcstsep Lachnospiraceae\dcstsep\taxon{Alistipes}
dominance-lineage at depth 3 (adjusted OR 10.15, q = 6.2e-22). The
depth-4
\taxon{Bacteroides}\dcstsep Lachnospiraceae\dcstsep\taxon{Alistipes}\dcstsep\taxon{Faecalibacterium}
dominance-lineage retained a finite Bayesian estimate but became too small for the
primary adjusted frequentist screen. The rare \taxon{Morganella}
and \taxon{Proteus} states likewise fell below the prespecified support
threshold in this sensitivity analysis, reinforcing the need to interpret
them cautiously (Supplementary Table S7).

Coarse antibiotic-status sensitivity analyses supported the same
calibration. Among age-, sex-, and BMI-complete samples, adding the AGP
\texttt{Antibiotic\_Status} field as a covariate or excluding samples
reporting antibiotic use in the previous week, month, or six months did
not attenuate the leading \taxon{Bacteroides}-centered IBD signals. The
depth-2 \taxon{Bacteroides}\dcstsep Lachnospiraceae dominance-lineage,
for example, had adjusted OR 3.21 in the age/sex/BMI model, 3.26 after
including coarse antibiotic status, and 3.44 after excluding recent
antibiotic users. These analyses do not replace detailed medication,
treatment, or disease-activity data, but they make recent antibiotic use
an unlikely sole explanation for the main AGP IBD result
(Supplementary Table S9). The selected IBD labels were also stable when
compared against the pure-policy dCST baseline, indicating that the
primary \taxon{Bacteroides}-centered findings are not artifacts of the
absorb low-frequency policy (Supplementary Table S12).

\begin{figure}[H]
\centering
\includegraphics[width=0.98\linewidth,height=\textheight,keepaspectratio,alt={IBD absorb dominance-lineage figure. Panel A shows the progression from the broad depth-1 Bacteroides state to deeper Bacteroides-centered IBD-associated dominance-lineages at depths 2 to 4. Panel B compares frequentist and Bayesian adjusted odds ratios for the main IBD dominance-lineage states and for rare enterobacterial states such as Morganella and Proteus.}]{../assets/figures/FIGURE_3_ibd_absorb_dominance_lineages.png}
\caption{IBD-associated dominance-lineages resolve across the absorb dCST
hierarchy. Panel A shows the progression from the broad depth-1
\taxon{Bacteroides} state to the depth-2 through depth-4
IBD-associated dominance-lineages emphasized in the main text. Panel B
compares frequentist and Bayesian adjusted odds ratios for
representative IBD-associated dominance-lineages and reports the corresponding
IBD/control counts. The common dominance-lineages are supported by both the
frequentist and Bayesian analyses. Rare enterobacterial dominance-lineages, including \taxon{Morganella}
and \taxon{Proteus}, are included as sparse-label examples: Bayesian
estimates remain finite, whereas the primary adjusted frequentist
screen skips labels with insufficient support.}
\end{figure}

External validation used the two complementary modes defined in Methods:
rebuilt-cohort analysis and direct transfer of AGP-derived absorb dCST
labels.

Across the completed inflammatory cohorts, rebuilt-cohort validation was
strongest in Halfvarson 2017, where the \texttt{IBD\_vs\_Healthy},
\texttt{Crohn\_vs\_Healthy}, and \texttt{UC\_vs\_Healthy} comparisons
all remained significant after multiple-testing correction. The leading
rebuilt labels differed across those contrasts but pointed to the same
broad inflammatory axis: depletion of \taxon{Segatella} sp. in pooled
IBD and Crohn disease together with enrichment of \taxon{Bacteroides}
\emph{dorei} in ulcerative colitis. PRJEB84421 provided narrower
support through
\texttt{OFG\_vs\_Healthy}, where a depth-2
\taxon{Faecalibacterium}\dcstsep\taxon{Bacteroides} dominance-lineage was depleted
(q = 0.019). By contrast, rebuilt \texttt{IBD\_vs\_Healthy} in HMP2 and
rebuilt \texttt{Crohn\_vs\_Healthy} in Gevers 2014 remained directional
but did not retain BH-significant absorb dCST dominance-lineages, which is not
surprising given repeated sampling, substantial clinical heterogeneity,
and the more specialized ecology of those cohorts.

Direct transfer of AGP-derived absorb dCST labels was stricter but still
informative. Halfvarson remained
clearly positive when samples were mapped directly onto AGP-defined
dominance-lineages: pooled IBD, Crohn, and UC comparisons all retained q < 0.001,
with the transferred signal centered on enriched \taxon{Bacteroides}
states and depleted \taxon{Faecalibacterium}-linked states. After
explicitly coding pooled HMP2 inflammatory bowel disease as one case
label, HMP2 also retained a weaker but significant depth-2 transferred
\taxon{Bacteroides}\dcstsep\taxon{Faecalibacterium} depletion in
\texttt{IBD\_vs\_Healthy} (q = 0.027). Gevers 2014 and PRJEB84421 both
mapped well into the AGP hierarchy but did not retain BH-significant
transferred states, and the current Jacobs IBS cohort was likewise
directional at most. We therefore interpret the present external data as
strongest support for IBD-associated dominance-pattern reproducibility,
with exact label portability clearly demonstrated in Halfvarson and
supported more selectively in pooled HMP2 (Figure 4; Supplementary Table
S4).

Because Halfvarson and HMP2 contain repeated samples, we also repeated
the leading validation summaries after retaining one detected sample per
subject. This sensitivity analysis narrowed the interpretation. The
Halfvarson sample-level transfer results remained strongest overall, but
the pooled IBD and Crohn transfer contrasts attenuated to suggestive
rather than corrected-significant subject-level evidence, whereas the UC
transfer signal remained significant (q = 0.0029). In HMP2, the
transferred depth-2 \taxon{Bacteroides}\dcstsep\taxon{Faecalibacterium}
signal remained significant in the subject-level sensitivity analysis
(q = 0.012). Direct transfer also showed phenotype-associated mapping
differences in Halfvarson, especially at deeper depths, where IBD cases
mapped more often than controls. We therefore treat direct AGP-derived
label transfer as strongest when both the association and the mapping
diagnostics are favorable, and we avoid interpreting nonmapping as
biological absence (Supplementary Tables S10 and S11).

\begin{table}[H]
\footnotesize
\caption{External validation summary for the main IBD-centered
dominance-state summaries. Rebuilt validation tests whether disease-linked
dominance structure is rediscovered when the hierarchy is fit within an
external cohort; direct AGP-derived label transfer tests whether the
discovery-cohort labels themselves remain informative when assigned
directly to external samples.}
\centering
\begin{tabularx}{\textwidth}{@{} l Y Y Y @{}}
\toprule
Cohort & Validation context & Evidence summary & Interpretation \\
\midrule
Halfvarson 2017 & Stool IBD, Crohn, and UC comparisons versus healthy controls; longitudinal sampling. & Sample-level rebuilt-cohort and AGP-derived label-transfer signals in all three inflammatory comparisons; subject-level sensitivity attenuated pooled IBD and Crohn transfer but retained UC transfer. & Strongest sample-level evidence, with subject-level checks supporting a calibrated label-portability claim. \\
HMP2 / IBDMDB & Pooled IBD versus healthy after explicit case recoding; mucosal gut-site samples with repeated longitudinal sampling. & Rebuilt mode was directional only, but AGP-derived label transfer retained a weaker depth-2 \taxon{Bacteroides}\dcstsep\taxon{Faecalibacterium} signal at sample level and after one-sample-per-subject sensitivity. & Selective transfer support, not a broad independent rediscovery. \\
PRJEB84421 & Orofacial granulomatosis versus healthy controls. & Rebuilt depth-2 \taxon{Faecalibacterium}\dcstsep\taxon{Bacteroides} depletion retained significance (q = 0.019), but AGP-derived label transfer did not. & Narrow inflammatory support without label-level portability. \\
Gevers 2014 & Pediatric stool Crohn disease versus healthy controls. & Samples mapped into the AGP hierarchy, but neither rebuilt-cohort nor AGP-derived label-transfer tests retained corrected absorb dCST dominance-lineage signal. & Informative mapping but not a positive replication. \\
Jacobs 2023 IBS 250bp & Baseline IBS versus healthy controls. & Directional at most, with no corrected AGP-derived label-transfer signal. & IBS validation remains inconclusive. \\
\bottomrule
\end{tabularx}
\end{table}

\begin{figure}[H]
\centering
\includegraphics[width=0.98\linewidth,height=\textheight,keepaspectratio,alt={External validation comparison figure. Panel A compares the strongest BH-adjusted sample-level signal recovered for each key cohort-comparison under rebuilt-cohort absorb validation and direct AGP-derived absorb dCST label transfer. Panel B shows AGP-derived hierarchy mapping coverage across depths 1 to 4 for the same comparisons. Halfvarson stays positive in both sample-level modes, HMP2 retains a weaker transferred depth-2 signal, and the other cohorts map into the hierarchy but lose exact-state significance when AGP-derived labels are assigned directly.}]{../assets/figures/FIGURE_4_external_validation_modes.png}
\caption{External validation separates reproducibility of the
IBD-associated dominance pattern from exact label portability. Panel A
compares the strongest BH-adjusted sample-level signal for each
cohort-comparison under rebuilt-cohort absorb validation and direct
AGP-derived absorb dCST label transfer. Panel B shows AGP-derived
hierarchy mapping coverage across depths 1 to 4 for the same
comparisons.}
\end{figure}

\subsection{Irritable bowel syndrome
(IBS)}\label{irritable-bowel-syndrome-ibs}


IBS was included as a secondary, hypothesis-generating gut phenotype
because its microbiome signal is typically subtler than in overt
inflammatory bowel disease
(\citeproc{ref-Pittayanon2019_IBS}{Pittayanon et al. 2019};
\citeproc{ref-Dinan2015_IBS}{Dinan and Cryan 2015};
\citeproc{ref-Cryan2012_gutbrain}{Cryan and Dinan 2012}). In the absorb dCST
structure, IBS showed a modest but internally consistent shift toward
enterobacterial dominance. At depth 1,
\taxon{Escherichia-Shigella} was enriched in IBS (adjusted OR 1.24, q =
0.002), while \taxon{Staphylococcus}, \taxon{Faecalibacterium}, and
\taxon{Prevotella\_9} were depleted. At depth 2, two
\taxon{Escherichia-Shigella}/\taxon{Bacteroides}-centered states
remained significant, and by depth 3 the clearest surviving frequentist
signal was an
\taxon{Escherichia-Shigella}/\taxon{Bacteroides}/Lachnospiraceae
dominance-lineage (adjusted OR 1.78, q = 0.018).

The Bayesian follow-up preserved the same ecological pattern at greater
depth, with strong posterior support for depletion of
\taxon{Faecalibacterium}- and \taxon{Prevotella\_9}-linked states and
for enrichment of deeper
\taxon{Escherichia-Shigella}/\taxon{Bacteroides}/Lachnospiraceae
dominance-lineages even when the frequentist q-values became sparse.
These secondary results are consistent with a milder dysbiosis-associated
dominance shift rather than evidence for a distinct IBS-specific
ecological program, in line with the current IBS
literature
(\citeproc{ref-Pittayanon2019_IBS}{Pittayanon et al. 2019};
\citeproc{ref-Tap2017_IBS}{Tap et al. 2017}).

\subsection{Autoimmune diseases}\label{autoimmune-diseases}


The AGP autoimmune phenotype is an umbrella label rather than a single
disease and therefore requires more cautious interpretation than IBD.
Even so,
autoimmune disease is an important category because the literature
repeatedly links loss of immune-regulatory commensals and expansion of
pathobionts to altered Treg/Th17 balance
(\citeproc{ref-Scher2013_Prevotella_RA}{Scher et al. 2013};
\citeproc{ref-Atarashi2013_Clostridia_Treg}{Atarashi et al. 2013};
\citeproc{ref-DeLuca2019_autoimmune}{De Luca and Shoenfeld 2019};
\citeproc{ref-Vatanen2018_T1D}{Vatanen et al. 2018}).

The broad AGP autoimmune phenotype provided a secondary, exploratory
immune-associated comparison. Its absorb pattern was compatible with
inflammation-associated dysbiosis but more heterogeneous than IBD. At depth 1, the dominant
signals were enrichment of \taxon{Bacteroides} (adjusted OR 1.40, q =
1.6e-12) together with depletion of \taxon{Prevotella\_9},
\taxon{Streptococcus}, \taxon{Faecalibacterium}, and
\taxon{Staphylococcus}. At depth 2, the hierarchy sharpened around
\taxon{Bacteroides}-centered dominance-lineages:
\taxon{Bacteroides}/Lachnospiraceae reached adjusted OR 2.15 (q =
5.1e-22), \taxon{Bacteroides}/[\emph{Eubacterium}]
\emph{eligens\_group} reached adjusted OR 21.95 (q = 3.2e-20), and
\taxon{Bacteroides}/\taxon{Alistipes} reached adjusted OR 1.70 (q =
2.5e-05). The same pattern continued deeper, with
\taxon{Bacteroides}/Lachnospiraceae/\taxon{Alistipes} reaching adjusted
OR 5.52 (q = 3.9e-21) at depth 3 and
\taxon{Bacteroides}/Lachnospiraceae/\taxon{Alistipes}/\taxon{Faecalibacterium}
remaining significant at depth 4 (adjusted OR 3.44, q = 6.1e-05).

The Bayesian analysis again helped mostly at the deeper levels. It
supported the same major \taxon{Bacteroides}-centered dominance-lineages as the
frequentist analysis, but it also retained support for additional depth-4
states such as
\taxon{Bacteroides}/\taxon{Akkermansia}/Lachnospiraceae/\taxon{Faecalibacterium}
and depleted
\taxon{Bacteroides}/\taxon{Faecalibacterium}/Lachnospiraceae/\taxon{Agathobacter}
states that no longer crossed the conventional q-value threshold. The
pattern is directionally similar to IBD but requires more cautious
interpretation: this umbrella phenotype appears enriched for an
inflammatory \taxon{Bacteroides}-centered ecology, yet the pooled
autoimmune label is too broad to support a disease-specific mechanistic
claim. These findings are best interpreted as evidence for a shared
inflammatory ecological axis rather than one causal autoimmune taxon
(\citeproc{ref-Scher2013_Prevotella_RA}{Scher et al. 2013};
\citeproc{ref-Vatanen2018_T1D}{Vatanen et al. 2018};
\citeproc{ref-Atarashi2013_Clostridia_Treg}{Atarashi et al. 2013};
\citeproc{ref-DeLuca2019_autoimmune}{De Luca and Shoenfeld 2019}).

\section{Materials and Methods}\label{methods}

The discovery analysis used AGP gut samples from the PRIME-derived SILVA
gut-project abundance table (\citeproc{ref-McDonald2018_AGP}{McDonald
et al. 2018}; \citeproc{ref-Quast2013_SILVA}{Quast et al. 2013}). We
restricted the analysis to the 34,711 AGP rows present in the PRIME
gut-project abundance table. Samples with total library size below
1,000 reads were excluded, removing 4,421 samples. Taxa were then
retained only if they reached a count of at least 2 in at least 5\% of
the retained samples, corresponding to a prevalence threshold of 1,515
samples in this cohort. This feature filter retained 274 taxa, and no
additional samples became all-zero after feature filtering, leaving
30,290 samples for the discovery analysis. Age was available for 27,372
samples (90.4\%), sex for 25,598 (84.5\%), and BMI for 26,541
(87.6\%), producing 23,198 complete cases for the absorb-policy dCST
dominance-lineage follow-up models. Because dCST construction was applied
after this feature-retention step, all dCST labels should be interpreted as
retained-feature dominance labels. They do not imply dominance over every
raw feature removed during preprocessing. Feature labels use the deepest
informative SILVA annotation available for each retained feature, so
labels can refer to species, genera, families, orders, or unresolved
features such as \emph{Unassigned}.

Phenotype flags were derived directly from the AGP \texttt{Phenotype}
metadata field by fixed substring matching. The five outcome variables
analyzed in the discovery screen were coded from the phrases
\texttt{irritable bowel syndrome},
\texttt{inflammatory bowel disease}, \texttt{autoimmune},
\texttt{acid reflux}, and \texttt{seasonal allergies}. For each outcome,
samples without the corresponding outcome-defining phrase served as
controls, and outcome flags were therefore not mutually exclusive. In the filtered cohort, case
counts were 3,105 for IBS, 1,166 for IBD, 2,872 for autoimmune disease,
2,859 for acid reflux, and 9,493 for seasonal allergies; corresponding
complete-case counts after requiring age, sex, and BMI were 2,551, 993,
2,426, 2,279, and 8,284, respectively (Supplementary Table S1). For the
AGP discovery cohort, IBD subtype, disease activity, medication exposure,
and detailed treatment histories were not available in a form that
supported clinical subtype modeling; the discovery IBD analyses are
therefore interpreted as self-reported phenotype associations rather than
as treatment-naive clinical subtype analyses. The coarser AGP
\texttt{Antibiotic\_Status} field was available for 22,939 of 23,198
age/sex/BMI complete-case samples and was used in sensitivity analyses
that either added antibiotic status as a categorical covariate or
excluded samples reporting antibiotic use in the previous week, month, or
six months (Supplementary Table S9).

Absorb-policy dCSTs were constructed from the within-sample dominance
order (Figure 1). A depth-1 state is defined by the dominant retained
taxon, and deeper states are defined by recursively appending the
second-ranked, third-ranked, and subsequent retained taxa. Because
assignment depends only on within-sample rank order, the label is
unchanged whether the input is expressed as raw counts, relative
abundances, or any other sample-wide rescaling that preserves the order
of taxa within each sample. We used the \texttt{absorb} low-frequency
policy throughout. At depth 1, only dominance patterns observed in at
least \(n_0 = 50\) samples were kept as named states. Samples in smaller
patterns were reassigned to the retained state with the largest abundance
among the retained taxa for that sample, with deterministic tie handling.
We then refined the hierarchy with \(n_0 = 25, 12, 10, 8,\) and \(5\)
for depths 2 through 6, respectively. At each depth, the same rule was
applied within the current parent state, producing ordered dCST
dominance-lineages.
Depths 1 through 4 are treated as the primary results in this manuscript,
with deeper levels used mainly to judge whether the hierarchy remained
structured.

The pure-policy dCST hierarchy without absorption served as a baseline
comparison and is summarized in the supplement. The inferential workflow
reported here is therefore absorb based from the cohort overview through
phenotype-specific follow-up.

To ask whether a phenotype was associated with the overall absorb dCST
dominance-lineage distribution at a given depth, we first applied Fisher's exact test to
the phenotype-by-state table for that depth. We screened panels with
10,000 Monte Carlo draws and reran the strongest panels with 100,000
draws to obtain finer p-value resolution. We summarize the magnitude of
these depth-level associations with Cramer's V. We then performed
label-level follow-up in two parallel ways for IBD, IBS, autoimmune
disease, acid reflux, and seasonal allergies. The conventional analysis
used Fisher tests and logistic regression adjusted for age, sex, and BMI.
The Bayesian analysis used two complementary models. For unadjusted
label-level contrasts we used a beta-binomial prevalence model with a
Jeffreys prior, \(\mathrm{Beta}(0.5, 0.5)\), summarized by the posterior
median odds ratio, 95\% credible interval, and posterior probability
that the odds ratio exceeded 1. For adjusted analyses we fit Bayesian
logistic regression models with \texttt{arm::bayesglm} using the same
covariates as the frequentist regression and the package's default
weakly informative prior. Those adjusted Bayesian fits were summarized
by the exponentiated coefficient estimate, an approximate 95\% credible
interval based on the posterior standard error, and \(\Pr(\beta > 0)\)
for enrichment. Multiple testing in the conventional analyses was
controlled with the Benjamini-Hochberg procedure
(\citeproc{ref-BenjaminiHochberg1995}{Benjamini and Hochberg 1995}).
We treat BH-adjusted frequentist results as the primary evidence for
screen-level discovery and use Bayesian estimates as complementary
stabilized summaries, especially for sparse labels where maximum
likelihood estimates can be unstable or prespecified support filters
exclude labels from the primary adjusted screen.

For the primary IBD analysis, we repeated the follow-up after excluding
non-IBD controls with the other AGP phenotype strings analyzed in this
study (IBS, autoimmune disease, acid reflux, or seasonal allergies).
Selected IBD labels were also compared under the absorb and pure-policy
hierarchies to check whether the main signals depended on absorption of
low-frequency dominance patterns (Supplementary Table S12). External validation
was evaluated in two complementary ways. In rebuilt-cohort validation,
each external matrix was filtered and analyzed independently with the
same absorb workflow used for the discovery cohort. In direct
AGP-derived label-transfer validation, external samples were mapped onto
the AGP-derived absorb dCST hierarchy after aligning each cohort's taxa to the AGP
retained-feature dictionary. The completed validation cohorts included
HMP2 / IBDMDB, Halfvarson 2017, Gevers 2014, PRJEB84421, and Jacobs 2023
IBS 250bp; the main text emphasizes the IBD-focused cohorts because they
provided the clearest external readout. For cohorts with repeated
subjects, we repeated leading validation tests after retaining one
detected sample per subject, and for direct AGP-derived label transfer
we summarized case/control mapping rates by depth (Supplementary Tables
S10 and S11). The full comparison
summary is given in Supplementary Table S4. All manuscript figures,
tables, and analysis scripts were generated from workflows built around
the \texttt{linf} R package, with project-specific runners maintained in
the companion \texttt{gut\_microbiome} repository.

\section{Discussion}\label{general-discussion}

This study does not identify a new IBD dysbiosis axis. Instead, it
represents established IBD-associated gut dysbiosis as dominance states
that are interpretable at the level of individual samples. The central
pattern is a broad \taxon{Bacteroides} background that resolves into
deeper \taxon{Bacteroides}-centered dominance-lineages, together with
sparse enterobacterial states that remain exploratory because of their
small support and the absence of AGP disease-activity metadata. These
labels should not be interpreted as evidence that dominant taxa alone
explain IBD. Rather, they provide a compact way to name recurrent
ecological contexts in which IBD cases are overrepresented.

The contribution of this application is therefore representational and
analytical. The deterministic dCST label provides a rank-based
sample-level state, the absorb hierarchy preserves named deeper
dominance-lineages instead of sending a growing fraction of samples into
a residual tail, and the Bayesian layer keeps sparse-label estimates
finite while leaving BH-adjusted frequentist results as the primary
discovery framework. The secondary IBS and autoimmune analyses provide
related gut-outcome context rather than co-equal discovery claims; both
were directionally compatible with published dysbiosis patterns but
weaker or more heterogeneous than IBD.

Several limitations remain important. External validation supports
IBD-focused dominance-pattern reproducibility more strongly than
universal label portability: the strongest evidence comes from
Halfvarson, HMP2 provides selective transferred support, and several
cohorts mainly demonstrate mapping coverage rather than corrected
label-level replication. Repeated-subject sensitivity and case/control
mapping balance further argue against treating the current AGP-derived
label dictionary as universally transferable. Broader stool-based
validation will be needed, especially for IBS and any non-IBD extensions
of this approach. AGP phenotypes are also
self-reported and cross-sectional, so
residual confounding, medication effects, antibiotic exposure, and
reverse causation remain plausible even though the leading IBD signals
were robust to the coarse antibiotic-status checks available here. The dCST labels are also
retained-feature labels constructed after prevalence filtering, and some
low-resolution or potential non-gut labels appear in the full hierarchy;
those states are retained for transparent accounting but require caution
as biological signals. These concerns are especially important for
supplementary non-IBD phenotypes. The current evidence nevertheless
supports absorb-policy dCSTs as a biologically interpretable analysis
layer and motivates broader testing of their generalizability in
independent cohorts.

\section{Conclusion}\label{conclusion}

Taken together, these results support dCSTs as a compact, interpretable
way to summarize established IBD-associated gut dysbiosis as
sample-level dominance structure. The current evidence is strongest for
IBD, more selective for cross-cohort label portability, and exploratory
for broader gut-outcome applications.

\section{Data and Code
Availability}\label{data-and-code-availability}

The manuscript source and paper assets are maintained in the public
\texttt{linf} repository: \url{https://github.com/pgajer/linf}. The
project-specific discovery and follow-up runners are maintained in the
public \texttt{gut\_microbiome} repository:
\url{https://github.com/pgajer/gut_microbiome}. The discovery analysis
used the American Gut Project stool cohort (\texttt{PRJEB11419}) through
PRIME-derived gut abundance tables. The main project scripts
are:

\begin{itemize}
\tightlist
\item \path{run_agp_absorb_depthscan.R}, which generated the absorb
dCST hierarchy;
\item \path{run_absorb_label_followup.R}, which generated the
phenotype-specific frequentist and Bayesian follow-up analyses for
absorb-policy dCST dominance-lineages;
\item \path{run_external_case_control_validation_absorb.R}, which
generated the rebuilt-cohort validation results;
\item \path{run_external_case_control_validation_frozen_agp.R}, which
generated the direct AGP-derived label-transfer validation results;
\item \path{export_prime_project_validation_inputs.py}, which
generated PRIME-derived HMP2 input exports.
\item \path{build_ibd_dominance_lineage_profile_table.py}, which generated the
retained-feature abundance profile table for the main IBD dCST
dominance-lineages.
\item \path{build_ibd_reviewer_sensitivity_assets.R}, which generated
the antibiotic-status, subject-level validation, mapping-balance,
absorb-policy robustness, and taxonomic-rank sensitivity tables added
for reviewer calibration.
\end{itemize}

Supplementary materials summarize phenotype definitions and sample
counts, the pure-versus-absorb depth-change summary, IBD sensitivity
analyses, retained-feature dominance-lineage profiles, selected reference analyses
retained outside the main text, and the cross-cohort validation and
mapping-balance comparisons.

\section{Funding}\label{funding}

\paperfundingnote

\section{Competing Interests}\label{competing-interests}

The authors declare no competing interests.

\section{Acknowledgments}\label{acknowledgments}

We thank the participants and investigators of the American Gut Project,
HMP2 / IBDMDB, PRJEB84421, Halfvarson 2017, and Gevers 2014 for making
these cohort resources available, and we thank the PRIME project context
that made the large-scale gut screen feasible.

\section*{References}\label{references}
\addcontentsline{toc}{section}{References}

\protect\phantomsection\label{refs}
\begin{CSLReferences}{1}{0}
\bibitem[\citeproctext]{ref-Arumugam2011_enterotypes}
Arumugam, Manimozhiyan, Jeroen Raes, Eric Pelletier, et al. 2011.
{``Enterotypes of the Human Gut Microbiome.''} \emph{Nature} 473:
174--80. \url{https://doi.org/10.1038/nature09944}.

\bibitem[\citeproctext]{ref-Atarashi2013_Clostridia_Treg}
Atarashi, Koji, Taira Tanoue, Keiji Oshima, et al. 2013. {``Treg
Induction by a Rationally Selected Mixture of {C}lostridia Strains from
the Human Microbiota.''} \emph{Nature} 500: 232--36.
\url{https://doi.org/10.1038/nature12331}.

\bibitem[\citeproctext]{ref-BenjaminiHochberg1995}
Benjamini, Yoav, and Yosef Hochberg. 1995. {``Controlling the False
Discovery Rate: A Practical and Powerful Approach to Multiple
Testing.''} \emph{Journal of the Royal Statistical Society Series B} 57:
289--300. \url{https://doi.org/10.1111/j.2517-6161.1995.tb02031.x}.

\bibitem[\citeproctext]{ref-Budden2017_gutlung}
Budden, Kurtis F, Shaan L Gellatly, David L A Wood, et al. 2017.
{``Emerging Pathogenic Links Between Microbiota and the Gut--Lung
Axis.''} \emph{Nature Reviews Microbiology} 15: 55--63.
\url{https://doi.org/10.1038/nrmicro.2016.142}.

\bibitem[\citeproctext]{ref-Cani2007_endotoxemia}
Cani, Patrice D, Jacques Amar, Miguel A Iglesias, et al. 2007.
{``Metabolic Endotoxemia Initiates Obesity and Insulin Resistance.''}
\emph{Diabetes} 56: 1761--72. \url{https://doi.org/10.2337/db06-1491}.

\bibitem[\citeproctext]{ref-Cryan2012_gutbrain}
Cryan, John F, and Timothy G Dinan. 2012. {``Mind-Altering
Microorganisms: The Impact of the Gut Microbiota on Brain and
Behaviour.''} \emph{Nature Reviews Neuroscience} 13: 701--12.
\url{https://doi.org/10.1038/nrn3346}.

\bibitem[\citeproctext]{ref-Costea2018_enterotypes}
Costea, Paul I, Falk Hildebrand, Manimozhiyan Arumugam, et al. 2018.
{``Enterotypes in the Landscape of Gut Microbial Community
Composition.''} \emph{Nature Microbiology} 3: 8--16.
\url{https://doi.org/10.1038/s41564-017-0072-8}.

\bibitem[\citeproctext]{ref-DeSouza2021_GERD}
D'Souza, S M, K Houston, L Keenan, B S Yoo, P J Parekh, and D A Johnson.
2021. {``Role of Microbial Dysbiosis in the Pathogenesis of Esophageal
Mucosal Disease: A Paradigm Shift from Acid to Bacteria?''} \emph{World
Journal of Gastroenterology} 27: 2054--72.
\url{https://doi.org/10.3748/wjg.v27.i18.2054}.

\bibitem[\citeproctext]{ref-DeLuca2019_autoimmune}
De Luca, F, and Y Shoenfeld. 2019. {``The Microbiome in Autoimmune
Diseases.''} \emph{Clinical and Experimental Immunology} 195: 74--85.
\url{https://doi.org/10.1111/cei.13158}.

\bibitem[\citeproctext]{ref-Deshpande2018_esophageal}
Deshpande, N P, S M Riordan, N Castaño-Rodríguez, M R Wilkins, and N O
Kaakoush. 2018. {``Signatures Within the Esophageal Microbiome Are
Associated with Host Genetics, Age, and Disease.''} \emph{Microbiome} 6:
227. \url{https://doi.org/10.1186/s40168-018-0611-4}.

\bibitem[\citeproctext]{ref-Dinan2015_IBS}
Dinan, Timothy G, and John F Cryan. 2015. {``Irritable Bowel Syndrome: A
Microbiome-Gut-Brain Axis Disorder?''} \emph{Neurogastroenterology and
Motility} 27: 445--55. \url{https://doi.org/10.1111/nmo.12520}.

\bibitem[\citeproctext]{ref-Everard2013_Akkermansia}
Everard, Amandine, Clara Belzer, Lucie Geurts, et al. 2013.
{``Cross-Talk Between \emph{{A}kkermansia Muciniphila} and Intestinal
Epithelium Controls Diet-Induced Obesity.''} \emph{Proceedings of the
National Academy of Sciences} 110: 9066--71.
\url{https://doi.org/10.1073/pnas.1219451110}.

\bibitem[\citeproctext]{ref-Forslund2015_metformin}
Forslund, Kristoffer, Falk Hildebrand, Trine Nielsen, et al. 2015.
{``Disentangling Type 2 Diabetes and Metformin Treatment Signatures in
the Human Gut Microbiota.''} \emph{Nature} 528: 262--66.
\url{https://doi.org/10.1038/nature15766}.

\bibitem[\citeproctext]{ref-Frank2007_IBD}
Frank, Daniel N, April L St Amand, Robert A Feldman, et al. 2007.
{``Molecular-Phylogenetic Characterization of Microbial Community
Imbalances in Human Inflammatory Bowel Diseases.''} \emph{Proceedings of
the National Academy of Sciences} 104: 13780--85.
\url{https://doi.org/10.1073/pnas.0706625104}.

\bibitem[\citeproctext]{ref-Gevers2014_CD}
Gevers, Dirk, Subra Kugathasan, Lee A Denson, et al. 2014. {``The
Treatment-Naive Microbiome in New-Onset {C}rohn's Disease.''} \emph{Cell
Host and Microbe} 15: 382--92.
\url{https://doi.org/10.1016/j.chom.2014.02.005}.

\bibitem[\citeproctext]{ref-Gorenshtein2025_migraine}
Gorenshtein, Alon, Kamel Shihada, Liron Leibovitch, Tom Liba, and Avner
Goren. 2025. {``The Association Between Migraine and Gut Microbiota: A
Systematic Review.''} \emph{Acta Neurologica Belgica} 125: 977--87.
\url{https://doi.org/10.1007/s13760-025-02779-y}.

\bibitem[\citeproctext]{ref-Horn2022_Prevotella_lung}
Horn, Kadi J, Melissa A Schopper, Zoe G Drigot, and Sarah E Clark. 2022.
{``Airway \emph{{P}revotella} Promote {TLR2}-Dependent Neutrophil
Activation and Rapid Clearance of \emph{{S}treptococcus Pneumoniae} from
the Lung.''} \emph{Nature Communications} 13: 3321.
\url{https://doi.org/10.1038/s41467-022-31074-0}.

\bibitem[\citeproctext]{ref-Hoskinson2023_allergy}
Hoskinson, C L, D L Y Dai, K L Del Bel, et al. 2023. {``Delayed Gut
Microbiota Maturation in the First Year of Life Is a Hallmark of
Pediatric Allergic Disease.''} \emph{Nature Communications} 14: 4785.
\url{https://doi.org/10.1038/s41467-023-40336-4}.

\bibitem[\citeproctext]{ref-Hua2016_AGP_allergy}
Hua, Xinmin, James J Goedert, Angela Pu, Guoqin Yu, and Jianxin Shi.
2016. {``Allergy Associations with the Adult Fecal Microbiota: Analysis
of the {A}merican {G}ut {P}roject.''} \emph{EBioMedicine} 3: 172--79.
\url{https://doi.org/10.1016/j.ebiom.2015.11.038}.

\bibitem[\citeproctext]{ref-Jie2017_ASCVD}
Jie, Zhuye, Huihua Xia, Shi-Long Zhong, et al. 2017. {``The Gut
Microbiome in Atherosclerotic Cardiovascular Disease.''} \emph{Nature
Communications} 8: 845.
\url{https://doi.org/10.1038/s41467-017-00900-1}.

\bibitem[\citeproctext]{ref-Kappeter2023_migraine_FMT}
Kappéter, Ágnes, Dávid Sipos, Adorján Varga, et al. 2023. {``Migraine as
a Disease Associated with Dysbiosis and Possible Therapy with Fecal
Microbiota Transplantation.''} \emph{Microorganisms} 11: 2083.
\url{https://doi.org/10.3390/microorganisms11082083}.

\bibitem[\citeproctext]{ref-Karlsson2013_T2D_women}
Karlsson, Fredrik H, Valentina Tremaroli, Intawat Nookaew, et al. 2013.
{``Gut Metagenomics in {E}uropean Women with Normal, Impaired and
Diabetic Glucose Control.''} \emph{Nature} 498: 99--103.
\url{https://doi.org/10.1038/nature12198}.

\bibitem[\citeproctext]{ref-Knights2014_rethinking_enterotypes}
Knights, Dan, Tonya L Ward, Christopher E McKinlay, et al. 2014.
{``Rethinking Enterotypes.''} \emph{Cell Host \& Microbe} 16: 433--37.
\url{https://doi.org/10.1016/j.chom.2014.09.013}.

\bibitem[\citeproctext]{ref-Koren2013_enterotypes}
Koren, Omry, Dan Knights, Antonio Gonzalez, et al. 2013. {``A Guide to
Enterotypes Across the Human Body: Meta-Analysis of Microbial Community
Structures in Human Microbiome Datasets.''} \emph{PLoS Computational
Biology} 9: e1002863.
\url{https://doi.org/10.1371/journal.pcbi.1002863}.

\bibitem[\citeproctext]{ref-Lacy2021_ACG_IBS}
Lacy, Brian E, Mark Pimentel, Darren M Brenner, William D Chey, Laurie A
Keefer, Millie D Long, and Baha Moshiree. 2021. {``{ACG} Clinical
Guideline: Management of Irritable Bowel Syndrome.''} \emph{American
Journal of Gastroenterology} 116: 17--44.
\url{https://doi.org/10.14309/ajg.0000000000001036}.

\bibitem[\citeproctext]{ref-LeChatelier2013_richness}
Le Chatelier, Emmanuelle, Trine Nielsen, Junjie Qin, et al. 2013.
{``Richness of Human Gut Microbiome Correlates with Metabolic
Markers.''} \emph{Nature} 500: 541--46.
\url{https://doi.org/10.1038/nature12506}.

\bibitem[\citeproctext]{ref-Ley2006_obesity}
Ley, Ruth E, Peter J Turnbaugh, Samuel Klein, and Jeffrey I Gordon.
2006. {``Human Gut Microbes Associated with Obesity.''} \emph{Nature}
444: 1022--23. \url{https://doi.org/10.1038/4441022a}.

\bibitem[\citeproctext]{ref-LloydPrice2019_IBD}
Lloyd-Price, Jason, Cesar Arze, Ashwin N Ananthakrishnan, et al. 2019.
{``Multi-Omics of the Gut Microbial Ecosystem in Inflammatory Bowel
Diseases.''} \emph{Nature} 569: 655--62.
\url{https://doi.org/10.1038/s41586-019-1237-9}.

\bibitem[\citeproctext]{ref-Maaser2019_ECCO_ESGAR_IBD_diagnostics}
Maaser, Christian, Andreas Sturm, Stephan R Vavricka, et al. 2019.
{``{ECCO-ESGAR} Guideline for Diagnostic Assessment in {IBD} Part 1:
Initial Diagnosis, Monitoring of Known {IBD}, Detection of
Complications.''} \emph{Journal of Crohn's and Colitis} 13: 144--64.
\url{https://doi.org/10.1093/ecco-jcc/jjy113}.

\bibitem[\citeproctext]{ref-Manichanh2006_CD}
Manichanh, Chaysip, Laurence Rigottier-Gois, Emmanuelle Bonnaud, et al.
2006. {``Reduced Diversity of Faecal Microbiota in {C}rohn's Disease
Revealed by a Metagenomic Approach.''} \emph{Gut} 55: 205--11.
\url{https://doi.org/10.1136/gut.2005.073817}.

\bibitem[\citeproctext]{ref-McDonald2018_AGP}
McDonald, Daniel, Embriette Hyde, Justine W Debelius, et al. 2018.
{``{A}merican {G}ut: An Open Platform for Citizen Science Microbiome
Research.''} \emph{mSystems} 3: e00031--18.
\url{https://doi.org/10.1128/mSystems.00031-18}.

\bibitem[\citeproctext]{ref-McDonald2024_GG2}
McDonald, Daniel, Yueyu Jiang, Metin Balaban, et al. 2024.
{``{G}reen{G}enes2 Unifies Microbial Data in a Single Reference Tree.''}
\emph{Nature Biotechnology} 42: 715--18.
\url{https://doi.org/10.1038/s41587-023-02009-5}.

\bibitem[\citeproctext]{ref-Pittayanon2019_IBS}
Pittayanon, Rapat, Jennifer T Lau, Yuhong Yuan, Grigorios I Leontiadis,
Frances Tse, Michael G Surette, and Paul Moayyedi. 2019. {``Gut
Microbiota in Patients with Irritable Bowel Syndrome: A Systematic
Review.''} \emph{Gastroenterology} 157: 97--108.
\url{https://doi.org/10.1053/j.gastro.2019.03.049}.

\bibitem[\citeproctext]{ref-Qin2012_T2D}
Qin, Junjie, Yingrui Li, Zhiming Cai, et al. 2012. {``A Metagenome-Wide
Association Study of Gut Microbiota in Type 2 Diabetes.''} \emph{Nature}
490: 55--60. \url{https://doi.org/10.1038/nature11450}.

\bibitem[\citeproctext]{ref-Quast2013_SILVA}
Quast, Christian, Elmar Pruesse, Pelin Yilmaz, et al. 2013. {``The
{SILVA} Ribosomal {RNA} Gene Database Project: Improved Data Processing
and Web-Based Tools.''} \emph{Nucleic Acids Research} 41: D590--96.
\url{https://doi.org/10.1093/nar/gks1219}.

\bibitem[\citeproctext]{ref-Sartor2008_IBD}
Sartor, R Balfour. 2008. {``Microbial Influences in Inflammatory Bowel
Diseases.''} \emph{Gastroenterology} 134: 577--94.
\url{https://doi.org/10.1053/j.gastro.2007.11.059}.

\bibitem[\citeproctext]{ref-Scher2013_Prevotella_RA}
Scher, Jose U, Andrew Sczesnak, Randy S Longman, et al. 2013.
{``Expansion of Intestinal \emph{{P}revotella Copri} Correlates with
Enhanced Susceptibility to Arthritis.''} \emph{eLife} 2: e01202.
\url{https://doi.org/10.7554/eLife.01202}.

\bibitem[\citeproctext]{ref-Schubert2014_CDI}
Schubert, Allison M, Matthew A M Rogers, Corrie Ring, et al. 2014.
{``Microbiome Data Distinguish Patients with \emph{{C}lostridium
Difficile} Infection and Non-\emph{{C}. Difficile}-Associated Diarrhea
from Healthy Controls.''} \emph{mBio} 5: e01021--14.
\url{https://doi.org/10.1128/mBio.01021-14}.

\bibitem[\citeproctext]{ref-Seekatz2014_FMT}
Seekatz, Anna M, Jorn Aas, Christine E Gessert, et al. 2014. {``Recovery
of the Gut Microbiome Following Fecal Microbiota Transplantation.''}
\emph{mBio} 5: e00893--14. \url{https://doi.org/10.1128/mBio.00893-14}.

\bibitem[\citeproctext]{ref-Sehgal2021_CDI}
Sehgal, Kanika, and Sahil Khanna. 2021. {``Gut Microbiome and
\emph{{C}lostridioides Difficile} Infection: A Closer Look at the
Microscopic Interface.''} \emph{Therapeutic Advances in
Gastroenterology} 14: 1756284821994736.
\url{https://doi.org/10.1177/1756284821994736}.

\bibitem[\citeproctext]{ref-Sokol2008_Fprausnitzii}
Sokol, Harry, Bernardo Pigneur, Laetitia Watterlot, et al. 2008.
{``\emph{{F}aecalibacterium Prausnitzii} Is an Anti-Inflammatory
Commensal Bacterium Identified by Gut Microbiota Analysis of {C}rohn
Disease Patients.''} \emph{Proceedings of the National Academy of
Sciences} 105: 16731--36. \url{https://doi.org/10.1073/pnas.0804812105}.

\bibitem[\citeproctext]{ref-Strachan1989_hygiene}
Strachan, D P. 1989. {``Hay Fever, Hygiene, and Household Size.''}
\emph{BMJ} 299: 1259--60.
\url{https://doi.org/10.1136/bmj.299.6710.1259}.

\bibitem[\citeproctext]{ref-Takeuchi2023_carbohydrate}
Takeuchi, Takeshi, Tetsuya Kubota, Yumiko Nakanishi, et al. 2023. {``Gut
Microbial Carbohydrate Metabolism Contributes to Insulin Resistance.''}
\emph{Nature} 621: 591--98.
\url{https://doi.org/10.1038/s41586-023-06466-x}.

\bibitem[\citeproctext]{ref-Tang2015_TMAO_CKD}
Tang, W H Wilson, Zeneng Wang, David J Kennedy, et al. 2015. {``Gut
Microbiota-Dependent Trimethylamine {N}-Oxide ({TMAO}) Pathway
Contributes to Both Development of Renal Insufficiency and Mortality
Risk in Chronic Kidney Disease.''} \emph{Circulation Research} 116:
448--55. \url{https://doi.org/10.1161/CIRCRESAHA.116.305360}.

\bibitem[\citeproctext]{ref-Tang2013_TMAO}
Tang, W H Wilson, Zeneng Wang, Bruce S Levison, et al. 2013.
{``Intestinal Microbial Metabolism of Phosphatidylcholine and
Cardiovascular Risk.''} \emph{New England Journal of Medicine} 368:
1575--84. \url{https://doi.org/10.1056/NEJMoa1109400}.

\bibitem[\citeproctext]{ref-Tap2017_IBS}
Tap, Julien, Muriel Derrien, Henrik Törnblom, Romain Brazeilles, Sonia
Cools-Portier, Joël Doré, Sanna Störsrud, Benjamin Le Nevé, and Magnus
Simrén. 2017. {``Identification of an Intestinal Microbiota Signature
Associated with Severity of Irritable Bowel Syndrome.''}
\emph{Gastroenterology} 152: 111--23.
\url{https://doi.org/10.1053/j.gastro.2016.09.049}.

\bibitem[\citeproctext]{ref-Trompette2014_fiber_asthma}
Trompette, Aurélien, Evan S Gollwitzer, Khyati Yadava, et al. 2014.
{``Gut Microbiota Metabolism of Dietary Fiber Influences Allergic Airway
Disease and Hematopoiesis.''} \emph{Nature Medicine} 20: 159--66.
\url{https://doi.org/10.1038/nm.3444}.

\bibitem[\citeproctext]{ref-Turnbaugh2006_energy}
Turnbaugh, Peter J, Ruth E Ley, Michael A Mahowald, Vincent Magrini,
Elaine R Mardis, and Jeffrey I Gordon. 2006. {``An Obesity-Associated
Gut Microbiome with Increased Capacity for Energy Harvest.''}
\emph{Nature} 444: 1027--31. \url{https://doi.org/10.1038/nature05414}.

\bibitem[\citeproctext]{ref-Vatanen2018_T1D}
Vatanen, Tommi, Eric A Franzosa, Randall Schwager, et al. 2018. {``The
Human Gut Microbiome in Early-Onset Type 1 Diabetes from the {TEDDY}
Study.''} \emph{Nature} 562: 589--94.
\url{https://doi.org/10.1038/s41586-018-0620-2}.

\bibitem[\citeproctext]{ref-Vaziri2013_CKD}
Vaziri, Nosratola D, Jakk Wong, Madeleine Pahl, et al. 2013. {``Chronic
Kidney Disease Alters Intestinal Microbial Flora.''} \emph{Kidney
International} 83: 308--15. \url{https://doi.org/10.1038/ki.2012.345}.

\bibitem[\citeproctext]{ref-Walters2014_meta_obesity}
Walters, William A, Zech Xu, and Rob Knight. 2014. {``Meta-Analyses of
Human Gut Microbes Associated with Obesity and {IBD}.''} \emph{FEBS
Letters} 588: 4223--33.
\url{https://doi.org/10.1016/j.febslet.2014.09.039}.

\bibitem[\citeproctext]{ref-Wang2024_GERD_MR}
Wang, K, S Wang, Y Chen, X Lu, D Wang, Y Zhang, W Pan, C Zhou, and D
Zou. 2024. {``Causal Relationship Between Gut Microbiota and Risk of
Gastroesophageal Reflux Disease: A Genetic Correlation and Bidirectional
{M}endelian Randomization Study.''} \emph{Frontiers in Immunology} 15:
1327503. \url{https://doi.org/10.3389/fimmu.2024.1327503}.

\bibitem[\citeproctext]{ref-Wang2011_TMAO}
Wang, Zeneng, Erin Klipfell, Brian J Bennett, et al. 2011. {``Gut Flora
Metabolism of Phosphatidylcholine Promotes Cardiovascular Disease.''}
\emph{Nature} 472: 57--63. \url{https://doi.org/10.1038/nature09922}.

\bibitem[\citeproctext]{ref-Zhu2016_TMAO_thrombosis}
Zhu, W, J C Gregory, E Org, et al. 2016. {``Gut Microbe-Generated
Trimethylamine {N}-Oxide from Dietary Choline Is Prothrombotic in
Subjects.''} \emph{Cell} 165: 111--24.
\url{https://doi.org/10.1016/j.cell.2016.02.011}.

\end{CSLReferences}

\end{document}
